% Options for packages loaded elsewhere
\PassOptionsToPackage{unicode}{hyperref}
\PassOptionsToPackage{hyphens}{url}
\documentclass[
  11pt,
]{article}
\usepackage{xcolor}
\usepackage[margin=1in]{geometry}
\usepackage{amsmath,amssymb}
\setcounter{secnumdepth}{-\maxdimen} % remove section numbering
\usepackage{iftex}
\ifPDFTeX
  \usepackage[T1]{fontenc}
  \usepackage[utf8]{inputenc}
  \usepackage{textcomp} % provide euro and other symbols
\else % if luatex or xetex
  \usepackage{unicode-math} % this also loads fontspec
  \defaultfontfeatures{Scale=MatchLowercase}
  \defaultfontfeatures[\rmfamily]{Ligatures=TeX,Scale=1}
\fi
\usepackage{lmodern}
\ifPDFTeX\else
  % xetex/luatex font selection
\fi
% Use upquote if available, for straight quotes in verbatim environments
\IfFileExists{upquote.sty}{\usepackage{upquote}}{}
\IfFileExists{microtype.sty}{% use microtype if available
  \usepackage[]{microtype}
  \UseMicrotypeSet[protrusion]{basicmath} % disable protrusion for tt fonts
}{}
\makeatletter
\@ifundefined{KOMAClassName}{% if non-KOMA class
  \IfFileExists{parskip.sty}{%
    \usepackage{parskip}
  }{% else
    \setlength{\parindent}{0pt}
    \setlength{\parskip}{6pt plus 2pt minus 1pt}}
}{% if KOMA class
  \KOMAoptions{parskip=half}}
\makeatother
\usepackage{graphicx}
\makeatletter
\newsavebox\pandoc@box
\newcommand*\pandocbounded[1]{% scales image to fit in text height/width
  \sbox\pandoc@box{#1}%
  \Gscale@div\@tempa{\textheight}{\dimexpr\ht\pandoc@box+\dp\pandoc@box\relax}%
  \Gscale@div\@tempb{\linewidth}{\wd\pandoc@box}%
  \ifdim\@tempb\p@<\@tempa\p@\let\@tempa\@tempb\fi% select the smaller of both
  \ifdim\@tempa\p@<\p@\scalebox{\@tempa}{\usebox\pandoc@box}%
  \else\usebox{\pandoc@box}%
  \fi%
}
% Set default figure placement to htbp
\def\fps@figure{htbp}
\makeatother
% definitions for citeproc citations
\NewDocumentCommand\citeproctext{}{}
\NewDocumentCommand\citeproc{mm}{%
  \begingroup\def\citeproctext{#2}\cite{#1}\endgroup}
\makeatletter
 % allow citations to break across lines
 \let\@cite@ofmt\@firstofone
 % avoid brackets around text for \cite:
 \def\@biblabel#1{}
 \def\@cite#1#2{{#1\if@tempswa , #2\fi}}
\makeatother
\newlength{\cslhangindent}
\setlength{\cslhangindent}{1.5em}
\newlength{\csllabelwidth}
\setlength{\csllabelwidth}{3em}
\newenvironment{CSLReferences}[2] % #1 hanging-indent, #2 entry-spacing
 {\begin{list}{}{%
  \setlength{\itemindent}{0pt}
  \setlength{\leftmargin}{0pt}
  \setlength{\parsep}{0pt}
  % turn on hanging indent if param 1 is 1
  \ifodd #1
   \setlength{\leftmargin}{\cslhangindent}
   \setlength{\itemindent}{-1\cslhangindent}
  \fi
  % set entry spacing
  \setlength{\itemsep}{#2\baselineskip}}}
 {\end{list}}
\usepackage{calc}
\newcommand{\CSLBlock}[1]{\hfill\break\parbox[t]{\linewidth}{\strut\ignorespaces#1\strut}}
\newcommand{\CSLLeftMargin}[1]{\parbox[t]{\csllabelwidth}{\strut#1\strut}}
\newcommand{\CSLRightInline}[1]{\parbox[t]{\linewidth - \csllabelwidth}{\strut#1\strut}}
\newcommand{\CSLIndent}[1]{\hspace{\cslhangindent}#1}
\setlength{\emergencystretch}{3em} % prevent overfull lines
\providecommand{\tightlist}{%
  \setlength{\itemsep}{0pt}\setlength{\parskip}{0pt}}
\usepackage{amsmath,amssymb,amsthm}
\setlength{\parindent}{0pt}
\setlength{\parskip}{0.5em}
\theoremstyle{definition}
\newtheorem{proposition}{Proposition}
\newtheorem{corollary}{Corollary}
\newtheorem{lemma}{Lemma}
\newcommand{\pospart}[1]{\left[#1\right]_{+}}
\usepackage{bookmark}
\IfFileExists{xurl.sty}{\usepackage{xurl}}{} % add URL line breaks if available
\urlstyle{same}
\hypersetup{
  pdftitle={A simple model of public status cues},
  hidelinks,
  pdfcreator={LaTeX via pandoc}}

\title{A simple model of public status cues}
\usepackage{etoolbox}
\makeatletter
\providecommand{\subtitle}[1]{% add subtitle to \maketitle
  \apptocmd{\@title}{\par {\large #1 \par}}{}{}
}
\makeatother
\subtitle{Short model with ideological weighting of gains and dependence
costs}
\author{}
\date{\vspace{-2.5em}21 September 2026}

\begin{document}
\maketitle

Let \(R\in\{0,1\}\) indicate whether a revisionist power occupies the
top position in the relevant trade hierarchy for a given country. Let
\(P\in\{0,1\}\) indicate whether that position has become publicly
visible, for example, through press coverage using explicit ranking
language.

I assume that policymakers' attention is sparse
(\citeproc{ref-gabaix2014}{Gabaix 2014}). This means that a change in
the new power's status competes for attention with other events in the
world and does not automatically produce a change in foreign policy. I
call the product \(q=RP\in\{0,1\}\) a political cue. It captures a
situation in which the emerging power not only attains its new rank but
also has that rank explicitly brought to public attention by the media
or public officials. Let salience \(s(q)\geq0\) satisfy \(s(1)>s(0)\).
In words, becoming the new number one and receiving public recognition
makes the emerging power more salient to policymakers, bringing it to
their attention. The function \(s(q)\) represents this attention as an
exogenous decision weight.

\section{Valence}\label{valence}

Let:

\begin{itemize}
\tightlist
\item
  \(G\geq0\): the baseline perceived gains from moving closer to the new
  great power;
\item
  \(D\geq0\): the baseline perceived costs of increased dependence on
  the emerging power;
\item
  \(\iota\in[-1,1]\): the leadership's ideological orientation toward
  the emerging power, where
\end{itemize}

\[
\begin{cases}
\iota>0 & \text{indicates a favorable view of the emerging power;}\\
\iota=0 & \text{indicates neutrality;}\\
\iota<0 & \text{indicates an unfavorable view of the emerging power.}
\end{cases}
\]

The quantities \(G\) and \(D\) are evaluated on a common utility scale
before ideological weighting. Ideology affects how much importance the
leadership assigns to each component. I define the net marginal value of
strengthening ties with China as

\[
\theta=(1+\iota)G-(1-\iota)D
      =(G-D)+\iota(G+D).
\]

An ideologically neutral leadership, \(\iota=0\), evaluates the
relationship according to \(G-D\). A favorable orientation increases the
weight on gains and reduces the weight on dependence costs; an
unfavorable orientation does the reverse. Both weights are nonnegative
and sum to two. At the endpoints, \(\iota=1\) places zero weight on
dependence costs and \(\iota=-1\) places zero weight on gains. Thus
ideology operates through the evaluation of the relationship: if
\(G=D=0\), then \(\theta=0\) for every ideological orientation.

Positive \(\theta\) makes movement toward China attractive before
adjustment costs; negative \(\theta\) favors hedging or distancing. When
\(G+D>0\), the evaluation equals zero at \(\iota_0=(D-G)/(G+D)\). For
example, with \(G=2\) and \(D=5\), an orientation of \(\iota=0.5\)
yields \(\theta=0.5\), whereas \(\iota=-0.5\) yields \(\theta=-6.5\).
Whether either evaluation produces an actual policy change also depends
on the costs and implementation conditions specified below. The
comparison between \(q=0\) and \(q=1\) holds \(G\), \(D\), and \(\iota\)
fixed.

\section{Players, actions, and
information}\label{players-actions-and-information}

There are two players:

\begin{enumerate}
\def\labelenumi{\arabic{enumi}.}
\tightlist
\item
  The political leadership, \(L\), chooses a policy directive
  \(d\in\mathbb{R}\), measured as a normalized adjustment relative to
  the existing policy. Positive values indicate movement toward China;
  negative values indicate hedging or distancing. A directive of \(d=0\)
  preserves the existing policy.
\item
  The diplomatic bureaucracy, \(B\), observes the directive and chooses
  \(e\in\{0,1\}\), where \(e=1\) means implementation.
\end{enumerate}

The realized policy is

\[
y=ed.
\]

The baseline model assumes complete information: \(q\), \(G\), \(D\),
\(\iota\), the salience and priority functions, and the costs are common
knowledge.

\section{Timing and payoffs}\label{timing-and-payoffs}

The game proceeds in the following order:

\begin{enumerate}
\def\labelenumi{\arabic{enumi}.}
\tightlist
\item
  The leadership chooses the directive \(d\).
\item
  The bureaucracy observes \(d\) and chooses \(e\).
\item
  The policy \(y=ed\) is realized, and payoffs are received.
\end{enumerate}

The leadership values policy changes that move the relationship in its
preferred direction, but must weigh that return against the costs of
departure from existing policy. Its payoff, normalized to zero when it
issues no directive, is

\[
u_L(d,e;q)
=
e\left[s(q)\theta d-\frac{d^2}{2}\right]
-c_L|d|,
\tag{1}
\]

The term \(s(q)\theta d\) captures the perceived return from the
direction and extent of an implemented adjustment. For positive
salience, movement toward China yields a positive return when
\(\theta>0\); when \(\theta<0\), a negative directive yields a positive
return by moving away from a relationship the leadership evaluates
unfavorably. Salience determines the weight this evaluation receives in
the policy decision. A public cue can therefore increase the incentive
to act in either direction while leaving the underlying evaluation
\(\theta\) unchanged.

The quadratic term \(d^2/2\) captures the increasing marginal cost of
larger implemented departures from existing policy, such as the
disruption involved in revising established commitments. Its curvature
is normalized to one for simplicity. This term limits the preferred
extent of adjustment even when the leadership strongly favors a
particular direction. The factor \(e\) multiplies both this cost and the
directional return because they arise only when the bureaucracy carries
out the directive.

By contrast, \(c_L|d|\), with \(c_L>0\), captures the political effort
and reputational exposure involved in formulating and defending a
directive, whether or not it is implemented. For example, proposing a
change in diplomatic voting instructions requires articulating and
defending the new position even if the bureaucracy subsequently leaves
voting practice unchanged. This cost rises with the size of the proposed
departure and applies to both directions. Conditional on implementation,
it creates a region of inaction: the leadership issues a nonzero
directive only when the initial marginal return \(s(q)|\theta|\) exceeds
\(c_L\).

The bureaucracy's payoff is

\[
u_B(e,d;q)
=
e\,\mathbf{1}\{d\neq0\}\,[K(q)-c_B],
\tag{2}
\]

where \(c_B>0\) is the organizational cost of changing routines and
\(K(q)\) is the value to the bureaucracy of implementing a clear and
coordinated priority. Assume that

\[
K(1)>K(0).
\]

The bureaucracy in the baseline model is directionally neutral. It
derives no utility from moving closer to China or to the United States.
Its decision depends on priority, clarity, and implementation costs.

\section{Equilibrium}\label{equilibrium}

The solution concept is subgame-perfect equilibrium, obtained by
backward induction.

\begin{lemma}[Bureaucratic response]
For any directive $d\neq0$, the bureaucracy implements if and only if $K(q)>c_B$. In the event of equality, the conservative tie-breaking rule $e=0$ applies.
\end{lemma}

\textit{Proof.} If \(d\neq0\), implementation yields \(K(q)-c_B\),
whereas non-implementation yields zero. Comparing the two directly gives
the condition. If \(d=0\), there is no policy to implement, and \(e=0\)
is adopted. \(\square\)

\begin{proposition}[Equilibrium with a directionally neutral cue]
The policy outcome in the subgame-perfect equilibrium is

\[
y^*(q)=
\mathbf{1}\{K(q)>c_B\}
\operatorname{sign}(\theta)
\left[s(q)|\theta|-c_L\right]_{+}.
\tag{3}
\]

where $\left[x\right]_{+}=\max\{x,0\}$. The leadership chooses $d^*(q)=y^*(q)$, anticipating the bureaucratic response specified in Lemma 1.

In particular, the cue can change whether a response occurs and how strong it is, but it does not determine its sign.
\end{proposition}

\textit{Proof.} If \(K(q)\leq c_B\), the bureaucracy does not implement
nonzero directives. Since the leadership pays \(c_L|d|\) even for a
directive that is not implemented, its optimal choice is \(d=0\), and
therefore \(y=0\).

\begin{minipage}{\linewidth}
If $K(q)>c_B$, every nonzero directive is implemented. The leadership solves

\[
\max_d\left\{s(q)\theta d-\frac{d^2}{2}-c_L|d|\right\}.
\]
\end{minipage}

For \(d>0\), the first-order condition is \(d=s(q)\theta-c_L\); this
solution is positive only when \(s(q)\theta>c_L\). For \(d<0\), the
condition is \(d=s(q)\theta+c_L\); it is negative only when
\(s(q)\theta<-c_L\). In all other cases, the subgradient condition is
satisfied at \(d=0\). Combining the three regions yields (3).
\(\square\)

\subsection{Immediate implications}\label{immediate-implications}

\begin{corollary}[The cue has no intrinsic valence]
Whenever $y^*(q)\neq0$,

\[
\operatorname{sign}\bigl(y^*(q)\bigr)=\operatorname{sign}(\theta).
\]

Thus, $\theta>0$ produces accommodation, whereas $\theta<0$ produces hedging.
\end{corollary}

\begin{corollary}[Activation through attention or coordination]
The cue activates a previously absent response when both post-cue conditions hold,

\[
K(1)>c_B
\quad\text{and}\quad
s(1)|\theta|>c_L,
\]

and at least one of them failed to hold before the cue:

\[
K(0)\leq c_B
\quad\text{or}\quad
s(0)|\theta|\leq c_L.
\]
\end{corollary}

\begin{corollary}[Benchmark with unchanged information and coordination]
If $s(1)=s(0)$ and $K(1)=K(0)$, the cue does not change the equilibrium policy. This is the benchmark in which the public label changes neither attention nor the capacity for coordination.
\end{corollary}

The general result therefore concerns magnitude:

\[
|y^*(1)|\geq|y^*(0)|
\]

when the cue weakly increases \(s\) and \(K\). The direction depends on
\(\theta\).

\section{Extensions}\label{extensions}

\subsection{Issue-specific costs}\label{issue-specific-costs}

To capture selective adjustment, let \(c_L=c_{Lk}\) and
\(\theta_k=(1+\iota)G_k-(1-\iota)D_k\) for each issue \(k\), where
\(G_k\) and \(D_k\) are the baseline gains and dependence costs relevant
to that issue. Conditional on implementation being available both before
and after the cue, the attention channel activates a previously absent
response on issues satisfying

\[
s(0)|\theta_k|\leq c_{Lk}<s(1)|\theta_k|.
\]

Issues with an already active response can also adjust further as
salience increases; issues whose costs remain above the post-cue
threshold remain unchanged. If the cue instead makes implementation
available for the first time, the activation conditions in Corollary 2
apply issue by issue. The composition of the response depends on the
distribution of valuations and costs and does not follow automatically
from an issue being classified as human rights or belonging to any other
substantive area.

\subsection{Persistence and incumbent
prominence}\label{persistence-and-incumbent-prominence}

Let \(\rho\) denote the expected persistence of the new hierarchy and
\(h\) the prominence of the displaced partner. A simple extension allows

\[
\frac{\partial s(1;\rho,h)}{\partial\rho}>0,
\qquad
\frac{\partial s(1;\rho,h)}{\partial h}>0,
\]

with analogous relationships for \(K\). A persistent ranking provides a
more stable basis for adaptation; displacing a prominent incumbent makes
the reclassification more visible. Neither condition determines the
valence of \(\theta\).

\section*{References}\label{references}
\addcontentsline{toc}{section}{References}

\protect\phantomsection\label{refs}
\begin{CSLReferences}{1}{0}
\bibitem[\citeproctext]{ref-gabaix2014}
Gabaix, Xavier. 2014. {``A Sparsity-Based Model of Bounded
Rationality.''} \emph{The Quarterly Journal of Economics} 129 (4):
1661--1710. \url{https://doi.org/10.1093/qje/qju024}.

\end{CSLReferences}

\end{document}
